\section{An exact conductor minimum over \texorpdfstring{$\Q(\sqrt{-23})$}{Q(sqrt(-23))}}\label{sec:example}

We now compute a small example completely. It demonstrates both a genuine extension obstruction and the cost of retaining a source condition.

Let
\[
 E=\Q(\sqrt{-23}),\qquad
 \theta=\frac{1+\sqrt{-23}}2,\qquad
 \theta^2-\theta+6=0.
\]
The reduced primitive forms of discriminant $-23$ are
\[
 (1,1,6),\quad(2,-1,3),\quad(2,1,3).
\]
The reduced-form inequality $a\le\sqrt{23/3}<3$ makes this list exhaustive. Hence $\Cl(E)\simeq C_3$, entirely in the minus part. The only roots of unity are $\pm1$.

The prime $\mathfrak q=(3,\theta-1)$ is nonprincipal. Indeed, the norm form is $u^2+uv+6v^2$; for $v\ne0$ it is at least $23v^2/4>3$, and for $v=0$ it is a square. Set
\[
 \beta=1+2\theta=2+\sqrt{-23},\qquad \N\beta=27.
\]
Its reduction is zero at $\mathfrak q$ and nonzero at the conjugate prime above $3$, so
\begin{equation}\label{eq:q3-cube}
                       \mathfrak q^3=(\beta).
\end{equation}
In particular the class of $\mathfrak q$ generates $\Cl(E)$.

\subsection{Why conductor 13 fails}

At $13$, the roots of $T^2-T+6$ are $5$ and $9$, so $13$ splits. The minus local $3$-group has order three. The ray exact sequence gives a group $B$ of order nine:
\[
 0\longrightarrow C_3\longrightarrow B\longrightarrow C_3\longrightarrow0.
\]
At the root $\theta=5$, the residue of $\beta/c(\beta)$ is $4$ modulo $13$, and
\[
                           4^4\equiv9\not\equiv1\pmod {13}.
\]
Thus its $3$-primary component is nontrivial. By \eqref{eq:q3-cube}, the triple of a minus lift of the class above $3$ is a nonzero local inertia element. The lift therefore has order nine and $B\simeq C_9$.

Its local inertia subgroup is $3C_9$. Every character $C_9\to\mu_3$ kills that subgroup. A prescribed nontrivial local cubic character therefore needs global order nine, even though its conductor norm remains $13^2=169$. This proves that the order increase in Theorem~\ref{thm:extension} can be necessary in an actual number field.

The splitting condition above $3$ is more restrictive. The minus Frobenius at $3$ generates $C_9$; requiring its value to be one kills the whole ray group. Thus \emph{no nontrivial character of any $3$-power order} on this fixed support satisfies that condition. Enlarging the allowed character order does not repair it.

\subsection{The first support which works}

\begin{theorem}[Exact minimum]\label{thm:minimum}
Among nontrivial finite $3$-power anticyclotomic characters of $E$ which are trivial on the decomposition groups above $3$ and ramify only at primes lying in split rational conjugate pairs, the minimum conductor norm is
\[
                              151^2=22801.
\]
Exactly two characters attain this minimum. Both have order three and conductor $(151)$.
\end{theorem}

\begin{proof}
For a split rational prime $q\ne3$, the minus local $3$-group is cyclic of order $3^{h_q}$, where $h_q=v_3(q-1)$. If $h_q=0$, no such character can ramify at $q$. Assume henceforth $h_q\ge1$.

Let $B_q$ be the minus $3$-ray group modulo $(q)$, and let $b$ be a minus lift of Frobenius above $3$. Its image generates the ordinary class group. Consequently $B_q$ is generated by its local group $R_q$ and $b$, with the single additional class relation
\[
                         3b=r_q\in R_q.
\]
Up to inversion and multiplication by a $3$-adic unit, $r_q$ is the $3$-primary residue of $\beta/c(\beta)$, by \eqref{eq:q3-cube}. Killing Frobenius therefore gives
\begin{equation}\label{eq:ray-quotient}
                 B_q/\langle b\rangle
                    \simeq R_q/\langle r_q\rangle.
\end{equation}
It is nontrivial exactly when $\beta/c(\beta)$ is a cube modulo $q$. Choosing the conjugate root in the residue field only inverts the ratio and does not change this test.

The exact finite calculation below includes every split prime $q\le151$ with $q\equiv1\pmod3$. The root $t$ is a root of $T^2-T+6$ and
\[
        u_q=\frac{1+2t}{3-2t}\pmod q,\qquad
        z_q=u_q^{(q-1)/3}\pmod q.
\]
\begin{center}
\begin{tabular}{rrrrrr}
\toprule
$q$ & $t$ & $u_q$ & $h_q$ & $z_q$ & $|B_q/\langle b\rangle|$\\
\midrule
13  &5 &4  &1&9 &1\\
31  &8 &13 &1&5 &1\\
73  &30&54 &2&64&1\\
127 &28&78 &2&19&1\\
139 &17&102&1&42&1\\
151 &43&79 &1&1 &3\\
\bottomrule
\end{tabular}
\end{center}
At $151$ the quotient is $C_3$, giving exactly two nontrivial characters. Their inertia is nontrivial, so the conductor is exactly $(151)$.

For minimality, unramified characters are killed by the condition at $3$ because the class of $\mathfrak q$ generates $\Cl(E)$. A character with one rational prime in its support and smaller conductor would occur in an earlier row and has just been excluded. A character with at least two distinct rational support primes has their product at least $13\cdot31=403>151$. Anticyclotomicity makes both conjugate primes ramify together. At residue characteristic $q\ne3$, a finite $3$-power character is trivial on the higher unit groups, which are pro-$q$; every ramified conductor exponent is one. Thus a multi-prime support has conductor norm at least $403^2>151^2$. This proves the claim without a squarefree-conductor assumption being added separately.
\end{proof}

The certificate uses only exact reduced-form enumeration, norms, modular roots, and modular exponentiation. The class-field interpretation is proved above, not delegated to floating-point computation. The file \code{checks/check_ray_characters.py} produces the full table and a JSON certificate.

\subsection{What this example repairs}

A two-prime alternative also illustrates the role of the support. At $13$ and $31$, both local groups are $C_3$ and both components of the relation $3b$ are nonzero. After changing their generators, killing $b$ imposes the relation $(1,1)$ on $C_3^2$. The two nontrivial compatible characters have local exponent vectors $(1,-1)$ and $(-1,1)$. Their conductor norm is $(13\cdot31)^2=162409$. The one-prime support $151$ is cheaper.

These constructions satisfy the elementary anticyclotomic, tame-ramification, and decomposition-group conditions. They do not verify the full upstream auxiliary-twist predicate. For example, for a primitive cubic value,
\[
                    v_3(1-\zeta_3)=\frac12.
\]
Summing over the two primes of a single split pair already gives one; over the two pairs above $13$ and $31$ it gives two. An arbitrarily small prescribed valuation budget is not met. Selmer and special-value conditions remain separate as well.

The example is a completed small arithmetic optimization problem and a rigorous test case for the transfer interface. It is not presented as a new solution of a broad conjecture or as a literature-priority claim.
